\documentclass{article}
\usepackage[utf8]{inputenc}
\usepackage[russian]{babel}
\usepackage{url}

\title{Лабораторная работа 2}

\author{Федотов Р.Э. Рычков Е.А.}
\date{Ноябрь 2021}

\begin{document}

\maketitle
\begin{abstract}
    Данный текст был взят с игрового сайта stopgame.ru. Мы не имеем никакого дохода с данной статьи, мы всего лишь использовали его для демонстрации функций \LaTeX 
    
    Поэтому настоящий автор данной статьи:
    \author{\textbf{Алексей Лихачев}}
\end{abstract}
\section*{Введение}

Даже более-менее реалистичную игру про гольф можно назвать головоломкой. Вы попадаете на поле, перед вами простирается большая территория, а где-то вдалеке расположена лунка. Совершите указанное число ударов — выступите хорошо. Закинете мяч в лунку за меньшее количество ходов — завоюете похвалу за выполненный игл или бёрди. A Little Golf Journey — головоломка в чистом виде, из которой убрали многие привычные элементы гольфа, предложив взамен уютную, обволакивающую атмосферу.

\section{Ничего лишнего}

В каком-то смысле игру можно назвать идеальной для тех, кто любит гольф, но не хочет разбираться во всех его тонкостях. Не всем охота узнавать разницу между айроном, драйвером и вудом — здесь менять клюшки вообще не нужно. Кому-то не нравится заниматься вычислениями на грине перед последним ударом — тут выполнять патты проще простого. Закрученных ударов и прочих подобных техник нет.

Даже самого гольфиста мы не видим — просто на поле падает мяч, а мы силой мысли его отбиваем в разные стороны. Есть обычные удары, а есть мощные: у первых можно регулировать направление полёта и силу удара, а у вторых — только направление. Локации разные (от мистического леса до пустыни с кактусами), но везде есть свои аналоги рафов и бункеров — из них совершить дальний удар не удаётся.


Проще пареной репы.
Больше правил нет. Может показаться, что игра очень лёгкая, и это действительно так —  но лишь для тех, кто хочет просто\footnote{Пример сноски} протащить мяч из одного угла карты в другой. Решение «загадки» обычно очевидно, и никаких трудностей за редким исключением не возникает. Гораздо интереснее после запуска уровня обратить внимание на правый верхний угол и увидеть, сколько ударов требуется для получения максимальной оценки. И тут A Little Golf Journey превращается в настоящую головоломку.

Начинаешь понимать, что дизайн полей продуман гораздо лучше, чем казалось. Попадая на довольно крупную карту и видя, что для четырёх звёзд достаточно лишь двух ударов, приходишь в недоумение — это же невозможно! Но невыполнимых условий нет, и решаешь всё получше изучить. Вот скала, от которой мяч может отскочить. Вот деревья, сквозь кроны которых шарик легко пролетает. Вот островок посреди озера — направление удара нужно определить корректно, иначе мяч упадёт в воду. Ну а грин вокруг лунки больше напоминает кратер — если мяч залтетит удачно, до цели он может докатиться сам.


Нет повести печальнее на свете…
Решать такого рода загадки увлекательно. Хотя зарабатывать много звёзд не требуется, внутренний перфекционист хочет получить наилучшую оценку. Но иногда всё же сдаёшься и откладываешь уровень до лучших времён — на некоторых полях достичь идеала очень трудно, даже если примерно понимаешь, что нужно делать. Игре отчаянно не хватает отмены последнего действия — конечно, перезапустить уровень не составляет труда, да и ударов обычно нужно не больше трёх, но каждый раз открывать меню и выбирать нужный пункт немного утомительно.

\section{Приятный отдых}
Даже когда охотишься за звёздами, атмосфера в A Little Golf Journey остаётся расслабляющей и умиротворяющей. Красиво сыгранные на пианино мелодии, простенький, но симпатичный визуальный стиль, доступный геймплей — гольф вообще спокойный вид спорта (если это не Mario Golf), но здесь всё особенно безмятежно. Здорово оформлена карта, по которой мы, как в Super Mario, передвигаемся от одного уровня к другому. Поначалу она бесцветная и тоскливая, но вокруг каждой пройденной лунки вырастают деревья и появляются яркие объекты.

Карта хороша ещё и тем, что с её помощью находишь секретные пути к бонусным уровням. Если видишь, что из какой-то точки явно может быть проложена ещё одна дорожка, возвращаешься к той лунке и изучаешь поле внимательнее. А там либо еле заметные ворота стоят, либо дерево странно покачивается — пробуешь отбить мяч в их сторону и открываешь бонусную лунку, до которой нужно добраться. Иногда натыкаешься на почти невидимые кубы — при взаимодействии с ними игра требует собрать за указанное время белые шарики.

A Little Golf Journey: Обзор
Для доступа к некоторым областям нужно набрать указанное количество звёзд и коллекционных предметов.
Есть и поводы возвращаться к старым полям — периодически на них появляются коллекционные предметы, которые называют просто «синими штуками». Поначалу собирать их легко: недолго думая отбиваешь мяч в их сторону, и всё. Но со временем связанные с ними пазлы становятся не менее заковыристыми, чем прохождение большой карты за два хода. В общем, секретов тут полно; A Little Golf Journey не требует искать каждую мелочь, зато награждает дополнительными развлечениями самых внимательных.

Попутно игра рассказывает простую историю, не используя для этого слишком много слов. Чем больше уровней вы проходите, тем чаще удаётся прочесть коротенькие письма, среди которых есть как абстрактные напутствия, так и переписка между некими X и Y. Кто такой X? Является ли Y игроком? Любовники ли X и Y или хорошие друзья? Отвечать на эти вопросы все будут по-разному, и такие письма отлично разбавляют геймплей, состоящий из ударов клюшкой по мячу.

A Little Golf Journey: Обзор
Поля не всегда выглядят нормально.
Я обожаю игры про гольф, поэтому A Little Golf Journey сразу привлекла моё внимание. О потраченном времени жалеть не пришлось, наоборот — в перерывах между другими играми я время от времени возвращаюсь в неё, чтобы пройти ещё пару уровней или попытаться завоевать побольше звёзд на уже пройденных полях. Если гольф нравится вам именно элементами пазла, а привычные для жанра тонкости вас не особо волнуют, то от этой игры вы наверняка получите большое удовольствие.

\textbf{Плюсы}: уютная атмосфера с приятной музыкой и очаровательным стилем; множество уровней, в том числе бонусных; ради секретов хочется возвращаться на посещённые поля — интересно, как пройти уровень за один или два удара.

\textbf{Минус}ы: нет кнопки отмены последнего действия — с ней игра была бы ещё более расслабляющей и умиротворяющей.

\begin{thebibliography}{Сайт}
\bibitem{1} Портал StopGame.ru: \url{ https://stopgame.ru/show/122912/a_little_golf_journey_review}
\end{thebibliography}

Введение и список литературы обычно не нумеруются. Каким образом убрать из заголовка раздела автоматически вставляемый командой section{} номер?

Нужно после команды до фигурных скобок добавить *

Что означает параметр команды bibitem, расположенный в квадратных скобках.

Нумерация списка

Как убрать дату из заголовка статьи?

Использовать  date{}, не заполняя фигурные скобки

\end{document}
